\chapter*{Abstract}
\addcontentsline{toc}{chapter}{Abstract}

% Separadores decimales y de millares en inglés dentro de este resumen.
% \decimalpoint (babel-spanish) evita que el punto decimal se convierta en coma en modo matemático.
{\decimalpoint\sisetup{output-decimal-marker={.}, group-separator={,}}

The AI4Mars dataset provides pixel-level segmentation masks, produced by crowdsourced
annotation, for tens of thousands of images taken on the surface of Mars. These masks have been
used almost exclusively to train neural networks. This work asks which quantitative indicators of
rock coverage and organisation can be derived from them reproducibly, and what the limits of
their validity are with respect to differences in annotation and to human judgement.

A classical image-processing procedure was implemented. For each mask it computes visible rock
coverage as the ratio of rock pixels to labelled pixels, and estimates the number of individual
rocks through connected components, the distance transform and a watershed guided by prominence
maxima, with explicit area and shape filters. It was applied to the \cNEscenas{} images from the
\textit{Curiosity} rover navigation camera that have a training mask.

Coverage could be computed for \cNCobCalc{} scenes (\cPctCobCalc); \cNCobPos{} (\cPctCobPos{}
of the total) had coverage greater than zero, with a median of \cCobMedVal{} over labelled
pixels, and no evidence was found that high values are explained by the labelled fraction of the
scene. The count, applied to the \cFlagOk{} eligible scenes, yielded \cNRocas{} rocks. The two
indicators show practically no linear association ($r = \cDepCCPearson$), although non-linear
measures detect a weak dependence. Against the \cNGold{} expert masks included in the dataset,
median coverage drops to \cGoldCobMedPos; a trained segmenter, used as a common instrument on
both samples, attributes most of that difference to the expert scenes being far less rocky, and
\cDescAnot{} percentage points to the annotation source. An exploratory evaluation with one
observer shows that the count systematically underestimates (weighted kappa \cVKappaPond): it
reproducibly counts the geometrically separable instances within the big-rock class, but not the
number of visible rocks, because that class covers only part of the rock and often groups
several blocks into a single polygon.

The result is therefore negative for counting: counting individual rocks from these masks is not
straightforward. It is, instead, diagnostic for the dataset, whose inconsistencies —expert scenes
of very different composition, differing annotation criteria, a boundary between big rock and
bedrock that is ambiguous even for experts, and partial annotation of blocks— are documented with
quantitative evidence. The work also delivers a results set with twenty-four per-image indicators
and an evaluation design that separates the effect of annotation from that of the images.

\vspace{0.4cm}
\noindent
\textbf{Keywords:} Mars; AI4Mars; semantic segmentation; watershed; mathematical morphology;
rock coverage; crowdsourced annotation; data quality; indicator validity.
}
